% Options for packages loaded elsewhere
\PassOptionsToPackage{unicode}{hyperref}
\PassOptionsToPackage{hyphens}{url}
%
\documentclass[
]{article}
\usepackage{lmodern}
\usepackage{amsmath}
\usepackage{ifxetex,ifluatex}
\ifnum 0\ifxetex 1\fi\ifluatex 1\fi=0 % if pdftex
  \usepackage[T1]{fontenc}
  \usepackage[utf8]{inputenc}
  \usepackage{textcomp} % provide euro and other symbols
  \usepackage{amssymb}
\else % if luatex or xetex
  \usepackage{unicode-math}
  \defaultfontfeatures{Scale=MatchLowercase}
  \defaultfontfeatures[\rmfamily]{Ligatures=TeX,Scale=1}
\fi
% Use upquote if available, for straight quotes in verbatim environments
\IfFileExists{upquote.sty}{\usepackage{upquote}}{}
\IfFileExists{microtype.sty}{% use microtype if available
  \usepackage[]{microtype}
  \UseMicrotypeSet[protrusion]{basicmath} % disable protrusion for tt fonts
}{}
\makeatletter
\@ifundefined{KOMAClassName}{% if non-KOMA class
  \IfFileExists{parskip.sty}{%
    \usepackage{parskip}
  }{% else
    \setlength{\parindent}{0pt}
    \setlength{\parskip}{6pt plus 2pt minus 1pt}}
}{% if KOMA class
  \KOMAoptions{parskip=half}}
\makeatother
\usepackage{xcolor}
\IfFileExists{xurl.sty}{\usepackage{xurl}}{} % add URL line breaks if available
\IfFileExists{bookmark.sty}{\usepackage{bookmark}}{\usepackage{hyperref}}
\hypersetup{
  hidelinks,
  pdfcreator={LaTeX via pandoc}}
\urlstyle{same} % disable monospaced font for URLs
\usepackage[margin=1in]{geometry}
\usepackage{color}
\usepackage{fancyvrb}
\newcommand{\VerbBar}{|}
\newcommand{\VERB}{\Verb[commandchars=\\\{\}]}
\DefineVerbatimEnvironment{Highlighting}{Verbatim}{commandchars=\\\{\}}
% Add ',fontsize=\small' for more characters per line
\usepackage{framed}
\definecolor{shadecolor}{RGB}{248,248,248}
\newenvironment{Shaded}{\begin{snugshade}}{\end{snugshade}}
\newcommand{\AlertTok}[1]{\textcolor[rgb]{0.94,0.16,0.16}{#1}}
\newcommand{\AnnotationTok}[1]{\textcolor[rgb]{0.56,0.35,0.01}{\textbf{\textit{#1}}}}
\newcommand{\AttributeTok}[1]{\textcolor[rgb]{0.77,0.63,0.00}{#1}}
\newcommand{\BaseNTok}[1]{\textcolor[rgb]{0.00,0.00,0.81}{#1}}
\newcommand{\BuiltInTok}[1]{#1}
\newcommand{\CharTok}[1]{\textcolor[rgb]{0.31,0.60,0.02}{#1}}
\newcommand{\CommentTok}[1]{\textcolor[rgb]{0.56,0.35,0.01}{\textit{#1}}}
\newcommand{\CommentVarTok}[1]{\textcolor[rgb]{0.56,0.35,0.01}{\textbf{\textit{#1}}}}
\newcommand{\ConstantTok}[1]{\textcolor[rgb]{0.00,0.00,0.00}{#1}}
\newcommand{\ControlFlowTok}[1]{\textcolor[rgb]{0.13,0.29,0.53}{\textbf{#1}}}
\newcommand{\DataTypeTok}[1]{\textcolor[rgb]{0.13,0.29,0.53}{#1}}
\newcommand{\DecValTok}[1]{\textcolor[rgb]{0.00,0.00,0.81}{#1}}
\newcommand{\DocumentationTok}[1]{\textcolor[rgb]{0.56,0.35,0.01}{\textbf{\textit{#1}}}}
\newcommand{\ErrorTok}[1]{\textcolor[rgb]{0.64,0.00,0.00}{\textbf{#1}}}
\newcommand{\ExtensionTok}[1]{#1}
\newcommand{\FloatTok}[1]{\textcolor[rgb]{0.00,0.00,0.81}{#1}}
\newcommand{\FunctionTok}[1]{\textcolor[rgb]{0.00,0.00,0.00}{#1}}
\newcommand{\ImportTok}[1]{#1}
\newcommand{\InformationTok}[1]{\textcolor[rgb]{0.56,0.35,0.01}{\textbf{\textit{#1}}}}
\newcommand{\KeywordTok}[1]{\textcolor[rgb]{0.13,0.29,0.53}{\textbf{#1}}}
\newcommand{\NormalTok}[1]{#1}
\newcommand{\OperatorTok}[1]{\textcolor[rgb]{0.81,0.36,0.00}{\textbf{#1}}}
\newcommand{\OtherTok}[1]{\textcolor[rgb]{0.56,0.35,0.01}{#1}}
\newcommand{\PreprocessorTok}[1]{\textcolor[rgb]{0.56,0.35,0.01}{\textit{#1}}}
\newcommand{\RegionMarkerTok}[1]{#1}
\newcommand{\SpecialCharTok}[1]{\textcolor[rgb]{0.00,0.00,0.00}{#1}}
\newcommand{\SpecialStringTok}[1]{\textcolor[rgb]{0.31,0.60,0.02}{#1}}
\newcommand{\StringTok}[1]{\textcolor[rgb]{0.31,0.60,0.02}{#1}}
\newcommand{\VariableTok}[1]{\textcolor[rgb]{0.00,0.00,0.00}{#1}}
\newcommand{\VerbatimStringTok}[1]{\textcolor[rgb]{0.31,0.60,0.02}{#1}}
\newcommand{\WarningTok}[1]{\textcolor[rgb]{0.56,0.35,0.01}{\textbf{\textit{#1}}}}
\usepackage{graphicx}
\makeatletter
\def\maxwidth{\ifdim\Gin@nat@width>\linewidth\linewidth\else\Gin@nat@width\fi}
\def\maxheight{\ifdim\Gin@nat@height>\textheight\textheight\else\Gin@nat@height\fi}
\makeatother
% Scale images if necessary, so that they will not overflow the page
% margins by default, and it is still possible to overwrite the defaults
% using explicit options in \includegraphics[width, height, ...]{}
\setkeys{Gin}{width=\maxwidth,height=\maxheight,keepaspectratio}
% Set default figure placement to htbp
\makeatletter
\def\fps@figure{htbp}
\makeatother
\setlength{\emergencystretch}{3em} % prevent overfull lines
\providecommand{\tightlist}{%
  \setlength{\itemsep}{0pt}\setlength{\parskip}{0pt}}
\setcounter{secnumdepth}{-\maxdimen} % remove section numbering
\ifluatex
  \usepackage{selnolig}  % disable illegal ligatures
\fi

\author{}
\date{\vspace{-2.5em}}

\begin{document}

\begin{center}\rule{0.5\linewidth}{0.5pt}\end{center}

title: ``Networks characterization from Birth Dead method with NetSim''
author: - Carlos Acosta output: bookdown::pdf\_document2: fig\_caption:
yes number\_sections: true urlcolor: blue toc: true fig\_caption: yes

header-includes: -

\usepackage{float}

\begin{center}\rule{0.5\linewidth}{0.5pt}\end{center}

\hypertarget{packages}{%
\section{Packages}\label{packages}}

\begin{Shaded}
\begin{Highlighting}[]
\FunctionTok{library}\NormalTok{(NetSim)  }\CommentTok{\# simulate bd time{-}tree networks: https://github.com/jjustison/NetSim}
\FunctionTok{library}\NormalTok{(TreeSim)  }\CommentTok{\# simulate bd time{-}trees}
\FunctionTok{library}\NormalTok{(ape)  }\CommentTok{\# tools for handling phylo objects}
\FunctionTok{library}\NormalTok{(diversitree)  }\CommentTok{\# tools for estimating lambda and mu}
\FunctionTok{library}\NormalTok{(phytools)}
\FunctionTok{library}\NormalTok{(}\StringTok{"dplyr"}\NormalTok{)  }\CommentTok{\# Data manipulation}
\FunctionTok{library}\NormalTok{(ggplot2)}
\FunctionTok{library}\NormalTok{(corrplot)}
\FunctionTok{library}\NormalTok{(cluster)}
\end{Highlighting}
\end{Shaded}

\hypertarget{using-netsim}{%
\section{Using NetSim:}\label{using-netsim}}

Simulation of 300 Trees with the following parameters using NetSim:

\begin{itemize}
\tightlist
\item
  lambda = 0.9
\item
  mu = 0.5
\item
  nu = 0.03
\item
  hybprops = c(0.5, 0.5, 0.5)
\end{itemize}

\begin{Shaded}
\begin{Highlighting}[]
\FunctionTok{set.seed}\NormalTok{(}\DecValTok{500}\NormalTok{)}
\NormalTok{numbsim }\OtherTok{=} \DecValTok{300}
\NormalTok{bdNS }\OtherTok{\textless{}{-}}\NormalTok{ NetSim}\SpecialCharTok{::}\FunctionTok{sim.bdh.taxa.ssa}\NormalTok{(}\AttributeTok{n =} \DecValTok{10}\NormalTok{, }\AttributeTok{numbsim =}\NormalTok{ numbsim, }\AttributeTok{lambda =} \FloatTok{0.9}\NormalTok{, }\AttributeTok{mu =} \FloatTok{0.5}\NormalTok{, }
    \AttributeTok{nu =} \FloatTok{0.03}\NormalTok{, }\AttributeTok{hybprops =} \FunctionTok{c}\NormalTok{(}\FloatTok{0.5}\NormalTok{, }\FloatTok{0.5}\NormalTok{, }\FloatTok{0.5}\NormalTok{), }\AttributeTok{hyb.inher.fxn =}\NormalTok{ NetSim}\SpecialCharTok{::}\FunctionTok{make.beta.draw}\NormalTok{(}\DecValTok{1}\NormalTok{, }
        \DecValTok{1}\NormalTok{), }\AttributeTok{frac =} \DecValTok{1}\NormalTok{, }\AttributeTok{mrca =} \ConstantTok{FALSE}\NormalTok{, }\AttributeTok{complete =} \ConstantTok{TRUE}\NormalTok{, }\AttributeTok{stochsampling =} \ConstantTok{FALSE}\NormalTok{, }\AttributeTok{hyb.rate.fxn =} \ConstantTok{NULL}\NormalTok{, }
    \AttributeTok{trait.model =} \ConstantTok{NULL}\NormalTok{)}
\end{Highlighting}
\end{Shaded}

As we can see in figure @ref(fig:chunk1), 59\% of the trees are extinct,
19.67\% are not reticulation and only 21.33\% are reticulated tree.

\begin{figure}
\centering
\includegraphics{NetDescription_files/figure-latex/chunk1-1.pdf}
\caption{(\#fig:chunk1)\label{fig:chunk1} Percentage of tree types using
NetSim with parameters lambda = 0.9, mu = 0.5 and nu = 0.03,}
\end{figure}

\hypertarget{extraction-of-information-from-tree}{%
\subsection{Extraction of information from
tree}\label{extraction-of-information-from-tree}}

We use the tree 35 (see fig.~@ref(fig:chunk4)) with three reticulations
to extract the tree information, and then extend of all the Trees. We
extract and analice the following information:

\begin{itemize}
\tightlist
\item
  Mean of edge length (meEdgeL)
\item
  cv of edge length (cvEdgeL)
\item
  Number of tips (TipNum)
\item
  Proportions of extintct taxa (PorExtin)
\item
  Number of internal nodes (NodeNum)
\item
  Time to reach 10 taxa (TimeTree)
\item
  Reticulation Number (RetNum)
\item
  Average time of birth of hybrids (TSHmean)
\item
  Minimum time of birth of hybrids (TSHmin)
\item
  Maximun time of birth of hybrids (TSHmax)
\end{itemize}

If reticulation number is one, TSHmean, TSHmin and TSHmax are the same,
if there are not reticulation the value is cero.

\begin{Shaded}
\begin{Highlighting}[]
\DocumentationTok{\#\# remove any trees with no taxa}
\NormalTok{bdNS }\OtherTok{\textless{}{-}}\NormalTok{ bdNS[}\SpecialCharTok{!}\FunctionTok{sapply}\NormalTok{(}\AttributeTok{X =}\NormalTok{ bdNS, }\AttributeTok{FUN =}\NormalTok{ is.null)]}

\CommentTok{\# Information about tree goes extinct=0 and no extinct tips are sampled=1}
\NormalTok{bdNS }\OtherTok{\textless{}{-}}\NormalTok{ bdNS[}\FunctionTok{sapply}\NormalTok{(}\AttributeTok{X =}\NormalTok{ bdNS, }\AttributeTok{FUN =}\NormalTok{ is.phylo)]  }\CommentTok{\#birth dead tree with NetSim with value 0 or 1}
\NormalTok{Tree2 }\OtherTok{\textless{}{-}}\NormalTok{ bdNS[[}\DecValTok{35}\NormalTok{]]  }\CommentTok{\#15\#10\#96}
\FunctionTok{plot}\NormalTok{(Tree2)}
\end{Highlighting}
\end{Shaded}

\begin{figure}
\centering
\includegraphics{NetDescription_files/figure-latex/chunk3-1.pdf}
\caption{(\#fig:chunk3)\label{fig:chunk4} Tree 35 which have three
reticulations}
\end{figure}

\begin{Shaded}
\begin{Highlighting}[]
\CommentTok{\# Tree2}
\end{Highlighting}
\end{Shaded}

\begin{Shaded}
\begin{Highlighting}[]
\DocumentationTok{\#\#\# Extracting general Information}
\NormalTok{meEdgeL }\OtherTok{\textless{}{-}} \FunctionTok{mean}\NormalTok{(Tree2}\SpecialCharTok{$}\NormalTok{edge.length)  }\CommentTok{\#mean of edge length}
\NormalTok{sdEdgeL }\OtherTok{\textless{}{-}} \FunctionTok{var}\NormalTok{(Tree2}\SpecialCharTok{$}\NormalTok{edge.length)}
\NormalTok{cvEdgeL }\OtherTok{\textless{}{-}}\NormalTok{ sdEdgeL}\SpecialCharTok{/}\NormalTok{meEdgeL  }\CommentTok{\#cv of edge length}
\NormalTok{TipNum }\OtherTok{\textless{}{-}} \FunctionTok{length}\NormalTok{(Tree2}\SpecialCharTok{$}\NormalTok{tip.label)  }\CommentTok{\# Number of tips}
\NormalTok{PorExtin }\OtherTok{\textless{}{-}}\NormalTok{ (TipNum }\SpecialCharTok{{-}} \DecValTok{10}\NormalTok{)}\SpecialCharTok{/}\NormalTok{TipNum  }\CommentTok{\#Proportions of extintct taxa }
\NormalTok{NodeNum }\OtherTok{\textless{}{-}}\NormalTok{ Tree2}\SpecialCharTok{$}\NormalTok{Nnode  }\CommentTok{\#Number of internal Nodes\# Maybe we can also divide it for TipNum}

\NormalTok{LT2 }\OtherTok{\textless{}{-}} \FunctionTok{ltt}\NormalTok{(Tree2, }\AttributeTok{plot =} \ConstantTok{FALSE}\NormalTok{)}
\NormalTok{Tim }\OtherTok{\textless{}{-}}\NormalTok{ LT2}\SpecialCharTok{$}\NormalTok{times}
\NormalTok{TimeTree }\OtherTok{\textless{}{-}} \FunctionTok{max}\NormalTok{(Tim)  }\CommentTok{\#Time to reach 10 taxa}

\CommentTok{\# extracting information about the Reticulations}

\NormalTok{Ret }\OtherTok{\textless{}{-}}\NormalTok{ Tree2}\SpecialCharTok{$}\NormalTok{reticulation}
\NormalTok{Ret1 }\OtherTok{\textless{}{-}}\NormalTok{ Ret}
\NormalTok{RetNum }\OtherTok{\textless{}{-}} \FunctionTok{nrow}\NormalTok{(Ret)  }\CommentTok{\# Reticulation Number}

\ControlFlowTok{if}\NormalTok{ (RetNum }\SpecialCharTok{\textgreater{}=} \DecValTok{1}\NormalTok{) \{}
    \ControlFlowTok{for}\NormalTok{ (i }\ControlFlowTok{in} \DecValTok{1}\SpecialCharTok{:}\FunctionTok{nrow}\NormalTok{(Ret1)) \{}
        \ControlFlowTok{for}\NormalTok{ (j }\ControlFlowTok{in} \DecValTok{1}\SpecialCharTok{:}\FunctionTok{ncol}\NormalTok{(Ret1)) \{}
\NormalTok{            Ret1[i, j] }\OtherTok{\textless{}{-}}\NormalTok{ Tim[}\FunctionTok{names}\NormalTok{(Tim) }\SpecialCharTok{\%in\%}\NormalTok{ Ret1[i, j]]}
\NormalTok{        \}}
\NormalTok{    \}}

\NormalTok{    TimeStartHibr }\OtherTok{\textless{}{-}} \FunctionTok{apply}\NormalTok{(Ret1, }\DecValTok{1}\NormalTok{, max)}\SpecialCharTok{/}\NormalTok{TimeTree  }\CommentTok{\#Time when birth the hibirds, we can divide it for the time to reach 10 taxa (TimeTree)}
\NormalTok{    TSHmean }\OtherTok{\textless{}{-}} \FunctionTok{mean}\NormalTok{(TimeStartHibr)  }\CommentTok{\# Average time of birth of hybrids }
\NormalTok{    TSHmin }\OtherTok{\textless{}{-}} \FunctionTok{min}\NormalTok{(TimeStartHibr)  }\CommentTok{\# minimum time of birth of hybrids }
\NormalTok{    TSHmax }\OtherTok{\textless{}{-}} \FunctionTok{max}\NormalTok{(TimeStartHibr)  }\CommentTok{\# maximun time of birth of hybrids }
\NormalTok{\} }\ControlFlowTok{else}\NormalTok{ \{}
\NormalTok{    TSHmean }\OtherTok{\textless{}{-}} \DecValTok{0}
\NormalTok{    TSHmin }\OtherTok{\textless{}{-}} \DecValTok{0}
\NormalTok{    TSHmax }\OtherTok{\textless{}{-}} \DecValTok{0}
\NormalTok{\}}

\CommentTok{\# If reticulation number is one, TSHmean, TSHmin and TSHmax are the same, if}
\CommentTok{\# there are not reticulation the value is cero.}
\NormalTok{dat }\OtherTok{\textless{}{-}} \FunctionTok{data.frame}\NormalTok{(TipNum, NodeNum, PorExtin, meEdgeL, cvEdgeL, TimeTree, RetNum, }
\NormalTok{    TSHmean, TSHmin, TSHmax)}

\NormalTok{dat}
\end{Highlighting}
\end{Shaded}

\begin{verbatim}
  TipNum NodeNum  PorExtin   meEdgeL   cvEdgeL TimeTree RetNum   TSHmean
1     14      20 0.2857143 0.3200865 0.4256297 1.995445      3 0.4342773
     TSHmin    TSHmax
1 0.2041249 0.6942674
\end{verbatim}

\hypertarget{extention-of-the-analisys-in-all-the-trees}{%
\subsection{Extention of the analisys in all the
trees}\label{extention-of-the-analisys-in-all-the-trees}}

We extended the analisys performed in section 2.1 on all trees

\begin{Shaded}
\begin{Highlighting}[]
\NormalTok{out1 }\OtherTok{\textless{}{-}} \FunctionTok{c}\NormalTok{()}
\ControlFlowTok{for}\NormalTok{ (k }\ControlFlowTok{in} \DecValTok{1}\SpecialCharTok{:}\FunctionTok{length}\NormalTok{(bdNS)) \{}
\NormalTok{    Tree2 }\OtherTok{\textless{}{-}}\NormalTok{ bdNS[[k]]}
    \DocumentationTok{\#\#\# Extracting general Information}
\NormalTok{    meEdgeL }\OtherTok{\textless{}{-}} \FunctionTok{mean}\NormalTok{(Tree2}\SpecialCharTok{$}\NormalTok{edge.length)  }\CommentTok{\#mean of edge length}
\NormalTok{    sdEdgeL }\OtherTok{\textless{}{-}} \FunctionTok{var}\NormalTok{(Tree2}\SpecialCharTok{$}\NormalTok{edge.length)}
\NormalTok{    cvEdgeL }\OtherTok{\textless{}{-}}\NormalTok{ sdEdgeL}\SpecialCharTok{/}\NormalTok{meEdgeL  }\CommentTok{\#cv of edge length}
\NormalTok{    TipNum }\OtherTok{\textless{}{-}} \FunctionTok{length}\NormalTok{(Tree2}\SpecialCharTok{$}\NormalTok{tip.label)  }\CommentTok{\# Number of tips}
\NormalTok{    PorExtin }\OtherTok{\textless{}{-}}\NormalTok{ (TipNum }\SpecialCharTok{{-}} \DecValTok{10}\NormalTok{)}\SpecialCharTok{/}\NormalTok{TipNum  }\CommentTok{\#Proportions of extintct taxa }
\NormalTok{    NodeNum }\OtherTok{\textless{}{-}}\NormalTok{ Tree2}\SpecialCharTok{$}\NormalTok{Nnode  }\CommentTok{\#Number of internal Nodes\# Maybe we can also divide it for TipNum}

\NormalTok{    LT2 }\OtherTok{\textless{}{-}} \FunctionTok{ltt}\NormalTok{(Tree2, }\AttributeTok{plot =} \ConstantTok{FALSE}\NormalTok{)}
\NormalTok{    Tim }\OtherTok{\textless{}{-}}\NormalTok{ LT2}\SpecialCharTok{$}\NormalTok{times}
\NormalTok{    TimeTree }\OtherTok{\textless{}{-}} \FunctionTok{max}\NormalTok{(Tim)  }\CommentTok{\#Time to reach 10 taxa}

    \CommentTok{\# extracting information about the Reticulations}

\NormalTok{    Ret }\OtherTok{\textless{}{-}}\NormalTok{ Tree2}\SpecialCharTok{$}\NormalTok{reticulation}
\NormalTok{    Ret1 }\OtherTok{\textless{}{-}}\NormalTok{ Ret}
\NormalTok{    RetNum }\OtherTok{\textless{}{-}} \FunctionTok{nrow}\NormalTok{(Ret)  }\CommentTok{\# Reticulation Number}

    \ControlFlowTok{if}\NormalTok{ (RetNum }\SpecialCharTok{\textgreater{}=} \DecValTok{1}\NormalTok{) \{}
        \ControlFlowTok{for}\NormalTok{ (i }\ControlFlowTok{in} \DecValTok{1}\SpecialCharTok{:}\FunctionTok{nrow}\NormalTok{(Ret1)) \{}
            \ControlFlowTok{for}\NormalTok{ (j }\ControlFlowTok{in} \DecValTok{1}\SpecialCharTok{:}\FunctionTok{ncol}\NormalTok{(Ret1)) \{}
\NormalTok{                Ret1[i, j] }\OtherTok{\textless{}{-}}\NormalTok{ Tim[}\FunctionTok{names}\NormalTok{(Tim) }\SpecialCharTok{\%in\%}\NormalTok{ Ret1[i, j]]}
\NormalTok{            \}}
\NormalTok{        \}}

\NormalTok{        TimeStartHibr }\OtherTok{\textless{}{-}} \FunctionTok{apply}\NormalTok{(Ret1, }\DecValTok{1}\NormalTok{, max)}\SpecialCharTok{/}\NormalTok{TimeTree  }\CommentTok{\#Time when birth the hibirds, we can divide it for the time to reach 10 taxa (TimeTree)}
\NormalTok{        TSHmean }\OtherTok{\textless{}{-}} \FunctionTok{mean}\NormalTok{(TimeStartHibr)  }\CommentTok{\# Average time of birth of hybrids }
\NormalTok{        TSHmin }\OtherTok{\textless{}{-}} \FunctionTok{min}\NormalTok{(TimeStartHibr)  }\CommentTok{\# minimum time of birth of hybrids }
\NormalTok{        TSHmax }\OtherTok{\textless{}{-}} \FunctionTok{max}\NormalTok{(TimeStartHibr)  }\CommentTok{\# maximun time of birth of hybrids }
\NormalTok{    \} }\ControlFlowTok{else}\NormalTok{ \{}
\NormalTok{        TSHmean }\OtherTok{\textless{}{-}} \DecValTok{0}
\NormalTok{        TSHmin }\OtherTok{\textless{}{-}} \DecValTok{0}
\NormalTok{        TSHmax }\OtherTok{\textless{}{-}} \DecValTok{0}
\NormalTok{    \}}

    \CommentTok{\# If reticulation number is one, TSHmean, TSHmin and TSHmax are the same, if}
    \CommentTok{\# there are not reticulation the value is cero.}
\NormalTok{    dat }\OtherTok{\textless{}{-}} \FunctionTok{data.frame}\NormalTok{(TipNum, NodeNum, PorExtin, meEdgeL, cvEdgeL, TimeTree, RetNum, }
\NormalTok{        TSHmean, TSHmin, TSHmax)}

\NormalTok{    out1 }\OtherTok{\textless{}{-}} \FunctionTok{rbind}\NormalTok{(out1, dat)}
\NormalTok{\}}

\FunctionTok{head}\NormalTok{(out1)}
\end{Highlighting}
\end{Shaded}

\begin{verbatim}
  TipNum NodeNum  PorExtin   meEdgeL   cvEdgeL TimeTree RetNum   TSHmean
1     15      15 0.3333333 0.5721684 0.9075774 3.317560      0 0.0000000
2     15      15 0.3333333 0.4118879 0.4096011 4.512402      0 0.0000000
3     14      14 0.2857143 0.8089697 0.8653747 6.793700      0 0.0000000
4     17      19 0.4117647 0.5390303 0.3657181 4.535813      1 0.8948177
5     19      25 0.4736842 0.4003979 0.2832645 4.542941      3 0.6990505
6     19      21 0.4736842 0.4611741 0.4221071 4.373698      1 0.6676135
     TSHmin    TSHmax
1 0.0000000 0.0000000
2 0.0000000 0.0000000
3 0.0000000 0.0000000
4 0.8948177 0.8948177
5 0.4606192 0.8503653
6 0.6676135 0.6676135
\end{verbatim}

\hypertarget{analisys-per-variable}{%
\subsection{Analisys per variable}\label{analisys-per-variable}}

In this section we analice the Trees varaible per variable (see fig
@ref(fig:chunk5)), as you can see, only PorExtin, meEdgeL, cvEdgeL, and
TimeTree have a distribution similar to the normal; NodeNum has
aHalf-normal distribution.

\begin{Shaded}
\begin{Highlighting}[]
\FunctionTok{par}\NormalTok{(}\AttributeTok{mfrow =} \FunctionTok{c}\NormalTok{(}\DecValTok{3}\NormalTok{, }\DecValTok{3}\NormalTok{))}
\NormalTok{nam }\OtherTok{\textless{}{-}} \FunctionTok{colnames}\NormalTok{(out1)}
\ControlFlowTok{for}\NormalTok{ (i }\ControlFlowTok{in} \FunctionTok{c}\NormalTok{(}\DecValTok{2}\SpecialCharTok{:}\DecValTok{6}\NormalTok{, }\DecValTok{8}\SpecialCharTok{:}\DecValTok{10}\NormalTok{)) \{}
    \FunctionTok{hist}\NormalTok{(out1[, i], }\AttributeTok{main =}\NormalTok{ nam[i])}
\NormalTok{\}}
\FunctionTok{barplot}\NormalTok{(}\FunctionTok{table}\NormalTok{(}\FunctionTok{sort}\NormalTok{(out1}\SpecialCharTok{$}\NormalTok{RetNum)), }\AttributeTok{main =} \StringTok{"Reticulation Number"}\NormalTok{)}
\end{Highlighting}
\end{Shaded}

\begin{figure}
\centering
\includegraphics{NetDescription_files/figure-latex/chunk63-1.pdf}
\caption{(\#fig:chunk63)\label{fig:chunk5} Univariate plot for all the
information extracted}
\end{figure}

\hypertarget{multivariate-analisys}{%
\subsection{Multivariate Analisys}\label{multivariate-analisys}}

In this part we perform a cor plot (see fig@ref(fig:chunk64)), in which
we can see that Reticulation number is positively correlated with
NodeNum. Besides we make a cluster analisys whit euclidean distance and
UPGMA grouped method(see fig @ref(fig:chunk7)), which sugest two cluster
numbers, so we can choose one to analyze with shaq.

\begin{figure}
\centering
\includegraphics{NetDescription_files/figure-latex/chunk73-1.pdf}
\caption{(\#fig:chunk73)\label{fig:chunk64} Corelation plot}
\end{figure}

\begin{Shaded}
\begin{Highlighting}[]
\NormalTok{clust }\OtherTok{\textless{}{-}} \FunctionTok{agnes}\NormalTok{(out1, }\AttributeTok{metric =} \StringTok{"euclidean"}\NormalTok{, }\AttributeTok{stand =} \ConstantTok{TRUE}\NormalTok{, }\AttributeTok{method =} \StringTok{"average"}\NormalTok{)}

\FunctionTok{par}\NormalTok{(}\AttributeTok{mfrow =} \FunctionTok{c}\NormalTok{(}\DecValTok{1}\NormalTok{, }\DecValTok{2}\NormalTok{))}
\FunctionTok{plot}\NormalTok{(clust)}
\end{Highlighting}
\end{Shaded}

\begin{figure}
\centering
\includegraphics{NetDescription_files/figure-latex/chunktr73-1.pdf}
\caption{(\#fig:chunktr73)\label{fig:chunk7} Cluster Analisys}
\end{figure}

\end{document}
